%! Characteristics of Exponential Functions — Unit 7, Lesson 7.0 — Guided Notes (Answer Key)
\documentclass[10pt]{article}
\usepackage{algebra2-article}
\usepackage{algebra2-key}   % pulls in -boxes; do NOT also load -boxes
\newcommand{\TallMath}[1]{$\displaystyle #1\rule[-1.4em]{0pt}{3.2em}$}
\newcommand{\lgrid}[4]{%
  \draw[step=1, linegray!40, very thin] (#1,#3) grid (#2,#4);
  \draw[->, semithick, charcoal] (#1,0)--(#2,0) node[below right, font=\scriptsize]{$x$};
  \draw[->, semithick, charcoal] (0,#3)--(0,#4) node[left, font=\scriptsize]{$y$};}
% Exponentials are plotted as a*e^{(ln b)x}: pgfmath's pow() is unreliable for negative and
% fractional exponents, but exp() is not. #2 is ln(b).
% ln2 = 0.693147   ln3 = 1.098612   ln4 = 1.386294
\newcommand{\expcurve}[4]{\draw[very thick, forest, domain=#3:#4, samples=120, smooth]
  plot (\x,{#1*exp((#2)*(\x))});}
% Vocab answer for the key: bold forest term, then the filled definition on a write-line.
% The leading and trailing \par are required: \ansline ends with \dotfill but does NOT end the
% paragraph, so without them each term label gets pulled onto the previous answer's line.
\newcommand{\vocabans}[2]{%
  \par\noindent\textbf{\textcolor{forest}{#1:}}\\[1pt]\ansline{#2}\par}

\begin{document}

\pageheader{Unit 7, Lesson 7.0}{Guided Notes --- Answer Key}
\namedateperiod

\vspace{0.08in}

\begin{objectivebox}
\textbf{By the end of this lesson, I will be able to\dots}
\begin{itemize}\setlength{\itemsep}{2pt}
  \item identify an \emph{exponential} function from its equation, its graph, or a \emph{table} (the \emph{constant ratio} fingerprint), and tell \emph{growth} from \emph{decay} by comparing the base to $1$;
  \item read the \emph{domain}, \emph{range}, \emph{intercepts}, \emph{increasing/decreasing} and \emph{positive/negative} intervals, \emph{extrema}, \emph{horizontal asymptote}, and \emph{end behavior} off an exponential graph;
  \item explain why an exponential graph has \emph{no} $x$-intercept and \emph{no} absolute minimum, even though it clearly has a floor.
\end{itemize}
\end{objectivebox}

\vspace{0.05in}

\begin{vocabbox}
\small
Fill in each term as we build it together.
\par\vspace{2pt}   % \par is required: \vocabans's \noindent is a no-op mid-paragraph
\vocabans{Exponential function}{a function with the variable in the \emph{exponent}, $f(x)=ab^{x}$, where $a\ne0$, $b>0$, and $b\ne1$}
\vocabans{Base $b$ (the growth/decay factor)}{the number being raised to the power --- what you multiply the output by each time $x$ increases by $1$}
\vocabans{Initial value $a$}{the output when $x=0$; since $b^{0}=1$, $f(0)=a$, so $a$ is always the $y$-intercept}
\vocabans{Exponential growth vs.\ exponential decay}{growth when $b>1$ (outputs increase); decay when $0<b<1$ (outputs shrink toward zero). Compare $b$ to $1$, not to $0$}
\vocabans{Constant ratio}{dividing each output by the one before it always gives the same number, $b$ --- the table fingerprint of an exponential function}
\vocabans{Horizontal asymptote (as a range boundary)}{a horizontal line the graph approaches but never reaches; for $y=ab^{x}$ it is $y=0$, and it is exactly the edge the range excludes}
\end{vocabbox}

\begin{hookbox}
Here are the two graphs this whole unit is built on: $y=2^{x}$ and $y=\left(\frac12\right)^{x}$.
\begin{center}
\begin{tikzpicture}[scale=0.36]
  % y = 2^x
  \begin{scope}
    \lgrid{-4}{4}{-1}{9}
    \begin{scope}
      \clip (-4,-1) rectangle (4,9);
      \expcurve{1}{0.693147}{-4}{3.17}
    \end{scope}
    \node[charcoal, font=\scriptsize] at (0,-2.6) {$y=2^{x}$};
  \end{scope}
  % y = (1/2)^x
  \begin{scope}[xshift=13cm]
    \lgrid{-4}{4}{-1}{9}
    \begin{scope}
      \clip (-4,-1) rectangle (4,9);
      \expcurve{1}{-0.693147}{-3.17}{4}
    \end{scope}
    \node[charcoal, font=\scriptsize] at (0,-2.6) {$y=\left(\tfrac12\right)^{x}$};
  \end{scope}
\end{tikzpicture}
\end{center}
\begin{itemize}\setlength{\itemsep}{1pt}
  \item Unit 6's square root graph gave up half the \emph{inputs} --- it simply stopped somewhere. These two graphs run all the way across, so the domain is back. What did they give up instead? \ansline{Half the \emph{outputs}. Both graphs stay entirely above the $x$-axis, so the range is only the positive numbers.}
  \item Neither graph ever crosses the $x$-axis. Warm-Up item~1 already told you why. What was it? \ansline{A power of a positive base is always positive --- a negative exponent gives a small \emph{fraction}, like $2^{-8}=\frac{1}{256}$, never $0$ and never negative.}
\end{itemize}
Radical functions made us ask \emph{where is the graph allowed to exist}. Exponential functions hand
the whole $x$-axis back and take something else away: \textbf{an entire half of the $y$-axis}, with a
line the graph gets closer and closer to but \emph{never} touches.
\end{hookbox}

\begin{notesbox}{1. What Makes a Function Exponential --- and How to Spot One in a Table}
Every family this year has kept the variable in the \textbf{base}. An \textbf{exponential function}
moves it into the \textbf{exponent}:
\[
  \underbrace{y=x^{2}}_{\text{\small variable is the \emph{base} --- quadratic}}
  \qquad\text{versus}\qquad
  \underbrace{y=2^{x}}_{\text{\small variable is the \emph{exponent} --- exponential}}
\]
The general form is
\[
  f(x) = a\,b^{\,x}, \qquad a\ne 0,\quad b>0,\quad b\ne 1 .
\]
\begin{itemize}[leftmargin=1.5em, itemsep=3pt]
  \item $a$ is the \textbf{initial value}. Because $b^{0}=$ \ans{1} for \emph{every} legal
        base, $f(0)=a$ --- so $a$ is always the \ans{$y$-intercept}.
  \item $b$ is the \textbf{base}, also called the growth or decay \textbf{factor}. It is what you
        \textbf{multiply} by every time $x$ goes up by $1$.
  \item Why must $b\ne 1$? Because $1^{x}=$ \ans{1} for every $x$, which would make the graph
        a \ans{horizontal} line --- not a new family at all.
  \item Why must $b>0$? A negative base collapses: $(-4)^{1/2}$ is not a real number, so the graph
        would be a scatter of holes instead of a smooth curve.
\end{itemize}

\vspace{3pt}
\textbf{The fingerprint test.} With evenly spaced $x$-values, a table names the family for you ---
no graph required.
\begin{center}
\renewcommand{\arraystretch}{1.35}
\begin{tabularx}{\linewidth}{@{}l l X@{}}
\hline
\textbf{Family} & \textbf{What stays constant} & \textbf{Example ($y$-values)} \\
\hline
Linear (Unit 2) & the \emph{difference} & $2,\ 5,\ 8,\ 11,\ 14$ \quad (add \ans{3} each time) \\
Quadratic (Unit 3) & the \emph{second} difference & $1,\ 4,\ 9,\ 16,\ 25$ \quad (differences $3,5,7,9$; those change by $+2$ each time) \\
\textbf{Exponential (Unit 7)} & the \textbf{ratio} & $2,\ 6,\ 18,\ 54,\ 162$ \quad (multiply by \ans{3} each time) \\
\hline
\end{tabularx}
\end{center}
That last row has a name: a \textbf{constant ratio}. Divide any output by the one before it ---
$\frac{6}{2}=$ \ans{3}, \; $\frac{18}{6}=$ \ans{3}, \; $\frac{54}{18}=$
\ans{3} --- and if you always get the same number, the function is exponential and that
number is \ans{$b$}.

\vspace{3pt}
\textbf{Growth or decay?} Compare the base to $1$ --- \emph{not} to $0$.
\begin{center}
\renewcommand{\arraystretch}{1.3}
\begin{tabularx}{\linewidth}{@{}l l X@{}}
\hline
$b>1$ & exponential \textbf{growth} & each step multiplies by more than $1$, so the outputs \ans{grow without bound} \\
$0<b<1$ & exponential \textbf{decay} & each step multiplies by a proper fraction, so the outputs \ans{shrink toward $0$} \\
\hline
\end{tabularx}
\end{center}
\end{notesbox}

\begin{notesbox}{2. The Growth Parent: $y=2^{x}$}
\begin{center}
\begin{minipage}[c]{0.38\linewidth}
\centering
\renewcommand{\arraystretch}{1.2}
\begin{tabular}{|c|c|}
\hline
$x$ & $y=2^{x}$ \\ \hline
$-3$ & \ans{$\frac18$} \\ \hline
$-1$ & \ans{$\frac12$} \\ \hline
$0$ & \ans{1} \\ \hline
$2$ & \ans{4} \\ \hline
$3$ & \ans{8} \\ \hline
\end{tabular}
\end{minipage}
\begin{minipage}[c]{0.58\linewidth}
\centering
\begin{tikzpicture}[scale=0.46]
  \lgrid{-4}{4}{-1}{9}
  \begin{scope}
    \clip (-4,-1) rectangle (4,9);
    \expcurve{1}{0.693147}{-4}{3.17}
  \end{scope}
  \draw[navy, dashed, semithick] (-4,0)--(4,0);
  \node[navy, font=\scriptsize, above left] at (-1.0,0) {$y=0$};
  \fill[charcoal] (0,1) circle (2.6pt);
  \fill[charcoal] (1,2) circle (2.6pt);
  \fill[charcoal] (2,4) circle (2.6pt);
  \fill[charcoal] (3,8) circle (2.6pt);
\end{tikzpicture}
\end{minipage}
\end{center}

\vspace{2pt}
\begin{enumerate}[leftmargin=1.7em, itemsep=3pt]
  \item \textbf{Domain:} \ans{$(-\infty,\infty)$} \quad \textbf{Range:} \ans{$(0,\infty)$} \\
        \emph{Watch the bracket.} $0$ is \textbf{not} in the range.
  \item \textbf{$y$-intercept:} \ans{$(0,1)$} \quad (that is $a=1$, exactly as predicted).
  \item \textbf{$x$-intercept / zeros:} \ans{none}. \; From Warm-Up item~1, $2^{x}$ is a
        \ans{positive} number for every real $x$, so it never equals $0$. \emph{This family never
        has a zero.}
  \item \textbf{Increasing / decreasing:} increasing on \ans{$(-\infty,\infty)$}; never
        \ans{decreasing}.
  \item \textbf{Positive / negative:} $y>0$ on \ans{$(-\infty,\infty)$}; \quad $y<0$
        \ans{nowhere}.
  \item \textbf{Extrema:} \ans{none}. There is no maximum (it climbs forever) --- and no minimum
        either, even though the graph clearly flattens out. Section 4 explains that second one.
  \item \textbf{Horizontal asymptote:} $y=$ \ans{0} \quad (the dashed line).
  \item \textbf{End behavior:} as $x\to+\infty$, $y\to$ \ans{$+\infty$} --- and \emph{fast}. \\
        As $x\to-\infty$, $y\to$ \ans{$0$}, flattening onto the asymptote from
        \ans{above}.
\end{enumerate}
\end{notesbox}

\begin{notesbox}{3. The Decay Parent: $y=\left(\frac12\right)^{x}$}
\begin{center}
\begin{minipage}[c]{0.38\linewidth}
\centering
\renewcommand{\arraystretch}{1.2}
\begin{tabular}{|c|c|}
\hline
$x$ & $y=\left(\frac12\right)^{x}$ \\ \hline
$-3$ & \ans{8} \\ \hline
$-1$ & \ans{2} \\ \hline
$0$ & \ans{1} \\ \hline
$2$ & \ans{$\frac14$} \\ \hline
$3$ & \ans{$\frac18$} \\ \hline
\end{tabular}
\end{minipage}
\begin{minipage}[c]{0.58\linewidth}
\centering
\begin{tikzpicture}[scale=0.46]
  \lgrid{-4}{4}{-1}{9}
  \begin{scope}
    \clip (-4,-1) rectangle (4,9);
    \expcurve{1}{-0.693147}{-3.17}{4}
  \end{scope}
  \draw[navy, dashed, semithick] (-4,0)--(4,0);
  \node[navy, font=\scriptsize, above right] at (1.0,0) {$y=0$};
  \fill[charcoal] (-3,8) circle (2.6pt);
  \fill[charcoal] (-2,4) circle (2.6pt);
  \fill[charcoal] (-1,2) circle (2.6pt);
  \fill[charcoal] (0,1) circle (2.6pt);
\end{tikzpicture}
\end{minipage}
\end{center}

\vspace{2pt}
The base is now a \emph{fraction}, so every step \textbf{shrinks} the output --- and the picture is
the growth parent's mirror image.
\begin{enumerate}[leftmargin=1.7em, itemsep=3pt]
  \item \textbf{Domain:} \ans{$(-\infty,\infty)$} \quad \textbf{Range:} \ans{$(0,\infty)$} \quad
        (\emph{identical} to the growth parent).
  \item \textbf{$y$-intercept:} \ans{$(0,1)$} \quad \textbf{$x$-intercept:} \ans{none}
  \item \textbf{Increasing / decreasing:} \ans{decreasing on $(-\infty,\infty)$}
  \item \textbf{Extrema:} \ans{none}, for the same reason as before.
  \item \textbf{Horizontal asymptote:} $y=$ \ans{0} --- the same line, but the graph now
        approaches it on the \ans{right} side instead of the left.
  \item \textbf{End behavior:} as $x\to+\infty$, $y\to$ \ans{$0$}; \; as $x\to-\infty$,
        $y\to$ \ans{$+\infty$}.
  \item \textbf{Why these two are secretly the same graph.} Since
        $\left(\frac12\right)^{x}=\dfrac{1}{2^{x}}$, the decay parent is the growth parent with the
        exponent \ans{negated ($2^{-x}$)}, and its graph is the growth graph reflected across the
        \ans{$y$}-axis. Lesson 7.2 makes this a rule.
\end{enumerate}
\end{notesbox}

\begin{notesbox}{4. The Horizontal Asymptote: a Floor the Graph Never Stands On}
A \textbf{horizontal asymptote} is a horizontal line the graph approaches as $x$ runs toward
$+\infty$ or $-\infty$. You met these in Unit 5 with rational functions. Here they take on a new job:
the asymptote is the \textbf{boundary of the range}.
\begin{center}
\renewcommand{\arraystretch}{1.2}
\begin{tabular}{|c|c|c|c|c|}
\hline
$x$ & $-1$ & $-4$ & $-10$ & $-20$ \\ \hline
$2^{x}$ & $0.5$ & $0.0625$ & $0.00098\ldots$ & $0.00000095\ldots$ \\ \hline
\end{tabular}
\end{center}
The outputs are being halved over and over. They get as close to $0$ as you like --- and they are
never \ans{$0$} and never \ans{negative}. Halving a positive number leaves a positive number,
forever.
\begin{itemize}[leftmargin=1.5em, itemsep=3pt]
  \item Therefore the range is $(0,\infty)$ and \textbf{not} $[0,\infty)$. Write the reason in your
        own words: \ansline{A square bracket would claim the function actually \emph{outputs} $0$ somewhere. It never does --- it only gets arbitrarily close --- so the endpoint is excluded.}
  \item Therefore there is \textbf{no absolute minimum}. A minimum must be a value the function
        \emph{actually reaches}, and $0$ is never reached.
\end{itemize}

\vspace{3pt}
\textbf{Compare this with Unit 6.} Both graphs below have a lowest edge. Only one of them
\emph{touches} it.
\begin{center}
\begin{tikzpicture}[scale=0.34]
  % sqrt(x): endpoint attained
  \begin{scope}
    \lgrid{-2}{8}{-2}{5}
    \begin{scope}
      \clip (-2,-2) rectangle (8,5);
      \draw[very thick, forest, domain=0:8, samples=90, smooth] plot (\x,{sqrt(\x)});
    \end{scope}
    \fill[forest] (0,0) circle (5pt);
    \node[charcoal, font=\scriptsize, align=center] at (3,-4.0)
      {$y=\sqrt{x}$ --- \emph{solid} endpoint:\\ the minimum $0$ \textbf{is} reached};
  \end{scope}
  % 2^x: asymptote never attained
  \begin{scope}[xshift=15cm]
    \lgrid{-4}{4}{-2}{5}
    \begin{scope}
      \clip (-4,-2) rectangle (4,5);
      \expcurve{1}{0.693147}{-4}{2.32}
    \end{scope}
    \draw[navy, dashed, semithick] (-4,0)--(4,0);
    \node[navy, font=\scriptsize, above left] at (-1.0,0) {$y=0$};
    \node[charcoal, font=\scriptsize, align=center] at (0,-4.0)
      {$y=2^{x}$ --- \emph{dashed} asymptote:\\ $0$ is \textbf{never} reached};
  \end{scope}
\end{tikzpicture}
\end{center}
\begin{itemize}[leftmargin=1.5em, itemsep=3pt]
  \item Unit 6's endpoint was \emph{in} the domain and \emph{in} the range, so it produced an
        absolute \ans{minimum}. An asymptote is in \ans{neither} of them.
  \item Unit 5's rational graphs had horizontal asymptotes too --- but a rational graph is
        \emph{allowed} to cross its horizontal asymptote in the middle. An exponential graph
        $y=ab^{x}$ never crosses its own, because $b^{x}$ is never \ans{$0$}.
  \item \textbf{The trade.} Unit 6 took away half the \ans{domain} (no negative radicands). Unit 7
        gives that back --- every real number is a legal exponent --- but takes away half the
        \ans{range} instead.
\end{itemize}
\end{notesbox}

\begin{notesbox}{5. The Full Feature List --- and How Growth and Decay Differ}
Read \emph{every} characteristic off the anchor $f(x)=4\left(\frac12\right)^{x}$.
\begin{center}
\begin{tikzpicture}[scale=0.44]
  \lgrid{-2}{6}{-1}{10}
  \begin{scope}
    \clip (-2,-1) rectangle (6,10);
    \expcurve{4}{-0.693147}{-1.32}{6}
  \end{scope}
  \draw[navy, dashed, semithick] (-2,0)--(6,0);
  \node[navy, font=\scriptsize, above right] at (3.8,0) {$y=0$};
  \fill[charcoal] (-1,8) circle (2.6pt) node[right, font=\scriptsize]{$(-1,8)$};
  \fill[charcoal] (0,4) circle (2.6pt) node[right, font=\scriptsize]{$(0,4)$};
  \fill[charcoal] (1,2) circle (2.6pt) node[above right, font=\scriptsize]{$(1,2)$};
  \fill[charcoal] (2,1) circle (2.6pt) node[above right, font=\scriptsize]{$(2,1)$};
\end{tikzpicture}
\end{center}
\begin{enumerate}[leftmargin=1.7em, itemsep=3pt]
  \item $a=$ \ans{4} and $b=$ \ans{$\frac12$}, so this is exponential \ans{decay}
        (because $b$ is \ans{less} than $1$).
  \item \textbf{Domain:} \ans{$(-\infty,\infty)$} \quad \textbf{Range:} \ans{$(0,\infty)$}
  \item \textbf{$y$-intercept:} \ans{$(0,4)$} \quad \textbf{$x$-intercept:} \ans{none}
  \item \textbf{Increasing / decreasing:} \ans{decreasing on $(-\infty,\infty)$}
  \item \textbf{Positive / negative:} $f(x)>0$ on \ans{$(-\infty,\infty)$}; \quad $f(x)<0$
        \ans{nowhere}.
  \item \textbf{Extrema:} \ans{none --- no absolute max and no absolute min}
  \item \textbf{Horizontal asymptote:} $y=$ \ans{0}, approached as $x\to$ \ans{$+\infty$}.
  \item \textbf{End behavior:} as $x\to+\infty$, $f(x)\to$ \ans{$0$}; \; as $x\to-\infty$,
        $f(x)\to$ \ans{$+\infty$}.
  \item \textbf{Check the constant ratio} using the labeled points: $\frac{f(1)}{f(0)}=\frac{2}{4}=$
        \ans{$\frac12$}, and $\frac{f(2)}{f(1)}=\frac{1}{2}=$ \ans{$\frac12$}. That is $b$, read
        straight off the graph.
\end{enumerate}

\vspace{4pt}
\textbf{Compare and contrast} the two exponential parents. Notice how many rows come out the
\emph{same} --- this family shares more than it differs.
\begin{center}
\renewcommand{\arraystretch}{1.3}
\begin{tabularx}{\linewidth}{@{}l X X@{}}
\hline
 & \textbf{Growth} $y=2^{x}$ & \textbf{Decay} $y=\left(\frac12\right)^{x}$ \\
\hline
Base $b$ & \ans{$2$} & \ans{$\frac12$} \\
Domain & \ans{$(-\infty,\infty)$} & \ans{$(-\infty,\infty)$} \\
Range & \ans{$(0,\infty)$} & \ans{$(0,\infty)$} \\
$y$-intercept & \ans{$(0,1)$} & \ans{$(0,1)$} \\
$x$-intercept & \ans{none} & \ans{none} \\
Increasing or decreasing & \ans{increasing} & \ans{decreasing} \\
Extrema & \ans{none} & \ans{none} \\
Horizontal asymptote & \ans{$y=0$} & \ans{$y=0$} \\
Asymptote approached on the\dots & \ans{left} & \ans{right} \\
\hline
\end{tabularx}
\end{center}
And against the families you already know: a \textbf{polynomial} could turn around many times; a
\textbf{rational} function broke apart at its vertical asymptotes; a \textbf{radical} function began
somewhere and stopped. An \textbf{exponential} function does none of those --- it is one smooth,
unbroken, never-turning curve living entirely on one side of a line it never crosses.
\end{notesbox}

\begin{practicebox}
The graph below is $p(x)=2\cdot 3^{x}$. Read each characteristic and write it on the line.
\begin{center}
\begin{tikzpicture}[scale=0.44]
  \lgrid{-3}{3}{-1}{10}
  \begin{scope}
    \clip (-3,-1) rectangle (3,10);
    \expcurve{2}{1.098612}{-3}{1.46}
  \end{scope}
  \draw[navy, dashed, semithick] (-3,0)--(3,0);
  \fill[charcoal] (0,2) circle (2.6pt) node[left, font=\scriptsize]{$(0,2)$};
  \fill[charcoal] (1,6) circle (2.6pt) node[right, font=\scriptsize]{$(1,6)$};
\end{tikzpicture}
\end{center}
\begin{enumerate}[leftmargin=1.7em, itemsep=3pt]
  \item $a=$ \ans{2} and $b=$ \ans{3}, so this is \ans{exponential growth}.
  \item \textbf{Domain:} \ans{$(-\infty,\infty)$} \quad \textbf{Range:} \ans{$(0,\infty)$}
  \item \textbf{$y$-intercept:} \ans{$(0,2)$} \quad \textbf{$x$-intercept:} \ans{none}
  \item \textbf{Horizontal asymptote:} \ans{$y=0$} \quad Approached as $x\to$ \ans{$-\infty$}.
  \item \textbf{End behavior:} as $x\to+\infty$, $p(x)\to$ \ans{$+\infty$}; \; as $x\to-\infty$,
        $p(x)\to$ \ans{$0$}.
  \item Does $p$ have an absolute maximum, an absolute minimum, both, or neither?
        \ans{neither} \\
        Explain in one sentence. \ansline{It climbs forever on the right, so there is no maximum;
        and on the left it only \emph{approaches} $y=0$ without ever reaching it, so there is no
        minimum either.}
\end{enumerate}
\end{practicebox}

\begin{teachernote}
\textbf{The one sentence to protect:} \emph{an asymptote is a boundary the range excludes, not a
value the function reaches.} Everything students get wrong today traces back to it --- the
$[0,\infty)$ range, the invented absolute minimum, and the hunt for an $x$-intercept are all the same
error wearing three costumes. The side-by-side picture in Section 4 (Unit 6's solid endpoint next to
Unit 7's dashed asymptote) is the single most important image of the lesson; do not rush it.

\textbf{Section 1 pacing.} Do not spend long on $b>0$ and $b\ne1$ --- they are guardrails, not
content, and the full algebraic story of fractional exponents is Lesson 6.1 material students
already own. The fingerprint table is what earns the time: it is the compare-and-contrast standard
(\textbf{A2.F.2b}, \textbf{A2.F.1e}) and it returns as the opening move of Lesson 7.5.

\textbf{Section 3, item 7} is a preview, not a requirement. Students only need to notice that
$\left(\frac12\right)^{x}=2^{-x}$; the reflection \emph{rule} is Lesson 7.2's job. If the class is
short on time, read it aloud and move on.

\textbf{Section 5, item 9} is worth cold-calling: reading $b$ off a \emph{graph} by dividing
consecutive outputs is exactly what Lesson 7.1 will ask students to do from a \emph{table}, and it
is the skill that makes 7.5's regression work interpretable.
\end{teachernote}

\end{document}
